\documentclass[11pt,a4paper]{article}

\usepackage[utf8]{inputenc}
\usepackage[T1]{fontenc}
\usepackage{geometry}
\usepackage{hyperref}
\usepackage{enumitem}
\usepackage{fancyhdr}
\usepackage{titlesec}

\geometry{margin=1in}
\hypersetup{colorlinks=true, linkcolor=blue, urlcolor=blue}

\pagestyle{fancy}
\fancyhf{}
\fancyhead[L]{\textbf{AICT Product Specification}}
\fancyhead[R]{\thepage}
\renewcommand{\headrulewidth}{0.4pt}

\title{\textbf{AICT Multi-Agent Management Platform} \\ \Large Product Specification}
\author{AICT Project}
\date{\today}

\begin{document}
\maketitle
\tableofcontents
\newpage

%% ============================================================
\section{Product Definition}
%% ============================================================

\subsection{One-line Pitch}

A web application where the \textbf{client} (user) talks to a \textbf{GM (General Manager)} agent who orchestrates \textbf{Operation Masters} and \textbf{Engineers} to build software---with a real-time kanban pipeline showing what every agent is doing, and a master specification folder (LaTeX to PDF) as the source of truth.

\subsection{Metaphor}

\begin{itemize}
  \item \textbf{Client} = the customer who hired a software firm.
  \item \textbf{GM} = the firm's general manager; single point of contact, reports to the client.
  \item \textbf{Operation Master} = project manager + QA; breaks down work, writes specs and acceptance tests, ensures engineers don't conflict.
  \item \textbf{Engineer} = developer; writes code, writes tests, ships.
\end{itemize}

The client never needs to talk to individual engineers. The client talks to the GM, sees the board, and reads the specs.

\subsection{Non-goals (Initial Release)}

\begin{itemize}
  \item No direct client-to-Engineer chat (always through GM or visible on the board).
  \item No CI/CD integration (agents run tests locally; CI is future).
  \item No multi-project support (one project/workspace at a time).
\end{itemize}

%% ============================================================
\section{Agent Roles}
%% ============================================================

\subsection{GM (exactly 1)}

\begin{itemize}
  \item \textbf{Owns}: Top-level specification, workflow definition, master spec folder (LaTeX sources and compiled PDFs).
  \item \textbf{Talks to}: Client (always), Operation Masters, Engineers.
  \item \textbf{Busy queue}: When GM is processing a client request or talking to another agent, incoming agent requests (from Operation Masters/Engineers) go into a FIFO queue. Client messages always take priority.
  \item \textbf{Responsibilities}: Understand client needs, translate to high-level specs, assign Operation Masters, approve major decisions, maintain the spec folder.
\end{itemize}

\subsection{Operation Master (0 or more, scaled by load)}

\begin{itemize}
  \item \textbf{Owns}: Requirement breakdown, spec-level acceptance tests, task assignment.
  \item \textbf{Talks to}: GM (queued if GM busy), Engineers under them, other Operation Masters.
  \item \textbf{Responsibilities}: Break GM specs into tasks, ensure tasks are non-conflicting (by module/feature/git branch), define ``done'' criteria (spec-level tests that must pass), review Engineer output against spec.
  \item \textbf{Git-aware}: Assigns work so Engineers operate on separate modules/branches to avoid merge conflicts.
\end{itemize}

\subsection{Engineer (multiple: Level 1 and Level 2)}

Two tiers by model and use:

\begin{itemize}
  \item \textbf{Engineer L1}: Well-balanced model (e.g.\ Claude 4.5). Used for routine implementation; good cost/performance. Owns implementation + unit/integration tests for assigned tasks.
  \item \textbf{Engineer L2}: Ultra-performance model (e.g.\ Claude 4.6). Used for hard or critical tasks; highest capability but token-heavy (``token nightmare''). Same responsibilities as L1, reserved for when L1 is insufficient or quality is paramount.
\end{itemize}

\noindent Both levels:
\begin{itemize}
  \item \textbf{Talks to}: GM (queued if GM busy), their Operation Master, other Engineers.
  \item \textbf{Responsibilities}: Implement assigned task, write tests, ensure tests pass, report completion.
\end{itemize}

%% ============================================================
\section{Client Surface}
%% ============================================================

The client interacts with the platform through three main views.

\subsection{View 1: Chat with GM}

\begin{itemize}
  \item Single persistent conversation thread between client and GM.
  \item Client describes what they want; GM responds with clarifications, plans, status updates.
  \item GM can attach or reference spec documents from the master spec folder.
  \item Shows GM status: \texttt{available} / \texttt{busy} (and what GM is busy with, e.g.\ ``talking to Operation Master-1'').
\end{itemize}

\subsection{View 2: Pipeline / Kanban Board (real-time)}

One unified board showing all work. Columns represent task status:

\begin{center}
  Backlog $\to$ Specifying $\to$ Assigned $\to$ In Progress $\to$ In Review $\to$ Done
\end{center}

\noindent Each card shows:
\begin{itemize}
  \item Task title and brief description.
  \item Assigned agent (Operation Master or Engineer) with role badge.
  \item Progress indicator (e.g.\ subtask completion, test pass rate).
  \item Current activity (e.g.\ ``Engineer-2 writing tests'', ``Operation Master-1 reviewing'').
  \item Git branch (if applicable).
  \item Time in current stage.
\end{itemize}

\noindent Additional board features:
\begin{itemize}
  \item \textbf{Progress bar}: Overall project completion (\% of tasks done).
  \item \textbf{Active task highlight}: Tasks currently being worked on are visually distinct.
  \item \textbf{Agent activity feed}: Small live feed showing what each agent is doing right now.
  \item \textbf{Filtering}: By agent, by status, by module/feature.
\end{itemize}

\subsection{View 3: Specification File Manager}

An online file management interface for the GM's master specification folder. The client can browse, organise, and edit specification files directly in the browser.

\subsubsection{Folder Browser}

\begin{itemize}
  \item \textbf{Tree / explorer view}: Hierarchical folder structure with expandable/collapsible directories, similar to a desktop file explorer or cloud drive.
  \item \textbf{Navigation}: Breadcrumb path bar at the top (e.g.\ \texttt{specs / backend / auth /}); click any segment to jump back.
  \item \textbf{File and folder icons}: Visual distinction between folders, \texttt{.tex} files, \texttt{.pdf} files, images, and other assets.
\end{itemize}

\subsubsection{File Operations}

\begin{itemize}
  \item \textbf{Create}: New folder, new \texttt{.tex} file (from template or blank), upload existing files.
  \item \textbf{Rename / Move}: Rename files or folders in place; drag-and-drop to move between folders.
  \item \textbf{Delete}: Delete files or folders (with confirmation).
  \item \textbf{Download}: Download any file or folder (as zip).
\end{itemize}

\subsubsection{In-browser Editor}

\begin{itemize}
  \item Click a \texttt{.tex} file to open it in an in-browser text editor with syntax highlighting.
  \item Auto-save or explicit save; changes are persisted to the server immediately.
  \item For \texttt{.tex} files: one-click \textbf{Compile} button triggers server-side \texttt{pdflatex}; the resulting PDF is shown in a side-by-side preview or opened in a new tab.
  \item For \texttt{.pdf} files: inline PDF viewer.
  \item For other text files: plain text editor.
\end{itemize}

\subsubsection{Client Authority and Versioning}

\begin{itemize}
  \item The client is the ultimate authority over the spec folder. Any edits the client makes override previous content.
  \item The GM is notified of client edits and adjusts the workflow accordingly.
  \item Each save is tracked with version history so the client and GM can review what changed and when.
\end{itemize}

%% ============================================================
\section{Guardrails}
%% ============================================================

\subsection{Git Safety}

\begin{itemize}
  \item \textbf{Blocked commands} (server-side): \texttt{git rebase}, \texttt{git reset --hard}, \texttt{git push --force}, \texttt{git commit --amend} (on pushed commits), \texttt{git filter-branch}, \texttt{git reflog expire}.
  \item \textbf{Implementation}: Backend wraps all git operations through a \texttt{GitService} that validates commands against a whitelist before execution.
  \item \textbf{Escalation}: If an agent requests a blocked command, the orchestrator pauses the task and sends an approval request to the client via the GM chat. Client approves or denies.
\end{itemize}

\subsection{Scope Isolation}

\begin{itemize}
  \item Engineers can only write to files within their assigned module/feature path.
  \item Operation Masters can only modify spec-related files and task definitions.
  \item GM can modify the spec folder and high-level project config.
\end{itemize}

\end{document}
